% Part I: provenance ----------------------------------------------------------------
\begin{table*}[!ht]
  \centering
  \footnotesize
  \setlength{\tabcolsep}{4pt}
  \caption{%
    Real-data task provenance.
    Most tasks are fetched by pinned OpenML data ID; the NSW data comes from
    the NBER Dehejia--Wahba mirror (screened CPS comparison sample). The split
    recreates the
    TableShift source and target pools locally without the TableShift runtime.
    The NSW task uses a raw outcome score (no source model).
    $N_{\mathrm{train}}$\,=\,source training rows;
    $N_{\mathrm{src}}$\,=\,held-out source evaluation rows;
    $N_{\mathrm{tgt}}$\,=\,target rows.%
  }
  \label{tab:appendix-provenance}
  \setlength{\tabcolsep}{3pt}
  \resizebox{\textwidth}{!}{%
  \begin{tabular}{l p{6em} l p{6em} p{5em} p{5em} r r r}
    \toprule
    Task         & OpenML dataset (ID)       & Domain          & Split variable        & Source pool           & Target pool            & $N_{\mathrm{train}}$ & $N_{\mathrm{src}}$ & $N_{\mathrm{tgt}}$ \\
    \midrule
    HELOC        & heloc (46932)             & Credit          & External Risk Estimate & $>63$                 & $\leq 63$              & 6{,}147 & 1{,}536 & 2{,}776 \\
    Readmission  & Diabetes130US (46922)     & Healthcare      & admission source id   & $\neq 7$              & $=7$                   & 26{,}583 & 6{,}645 & 4{,}000 \\
    Income       & ACSIncome (43141)         & Socioeconomic   & DIVISION              & $\neq$ New England    & $=$ New England        & 30{,}000 & 4{,}000 & 4{,}000 \\
    Coverage     & ACSPublicCoverage (43140) & Socioeconomic   & DIS                   & no disability         & disability ($=1$)      & 30{,}000 & 4{,}000 & 4{,}000 \\
    NSW Employment & n/a (NBER mirror)         & Employment      & treat                 & CPS comparison ($=0$)  & NSW treated ($=1$)     & 0        & 85      & 185 \\
    \bottomrule
  \end{tabular}%
  }
\end{table*}

% Part II: per-workflow results -----------------------------------------------------
\begin{table*}[!ht]
  \centering
  \footnotesize
  \setlength{\tabcolsep}{3.5pt}
  \caption{%
    Per-workflow $p$-values, decision flips, and effective sample sizes
    for the real-data tasks.
    \emph{Unw.}\,$p$\,=\,unweighted p-value;
    \emph{DW}\,$p$\,=\,doubly weighted p-value ($\lambda = 0.5$, 9{,}999 permutations).
    Flip\,=\,decision change at $\alpha = 0.05$ between unweighted and
    doubly weighted, in either direction (on HELOC Risk and Readmission rows
    the unweighted test rejects and the doubly weighted test does not; on
    HELOC Confidence and Income Risk rows only the doubly weighted test
    rejects); $\dagger$ marks a borderline flip.
    $N_{\mathrm{src}}$ and $N_{\mathrm{tgt}}$ are the original sample sizes;
    $\mathrm{ESS}_{\mathrm{src}}$ and $\mathrm{ESS}_{\mathrm{tgt}}$ are Kish's
    effective sample sizes for the doubly weighted mode.
    ESS is constant within each task because the domain classifier does not
    depend on the monitoring score.
    Reading a row---e.g., HELOC Risk: $N_{\mathrm{src}}=1{,}536$,
    $N_{\mathrm{tgt}}=2{,}776$, $\mathrm{ESS}=310$ and $337$---shows that
    most observations fall outside the common support, explaining why the
    unweighted rejection ($p=0.0001$) does not survive common-support
    restriction ($p=0.992$).
    Direction: $\uparrow$\,=\,higher target score is adverse;
    $\downarrow$\,=\,lower target score is adverse (score is higher-is-better).%
  }
  \label{tab:appendix-results}
  \begin{tabular}{l l c r r r r c r r}
    \toprule
    Task & Workflow & Dir. & $N_{\mathrm{src}}$ & $N_{\mathrm{tgt}}$ &
    Unw.~$p$ & DW~$p$ & Flip & $\mathrm{ESS}_{\mathrm{src}}$ & $\mathrm{ESS}_{\mathrm{tgt}}$ \\
    \midrule
    HELOC        & Risk       & $\uparrow$   & 1{,}536 & 2{,}776 & 0.0001 & 0.992 & $\checkmark$ & 310   & 337   \\
                 & Confidence & $\downarrow$ & 1{,}536 & 2{,}776 & 1.000 & 0.022 & $\checkmark$ & 310   & 337   \\
                 & Error      & $\uparrow$   & 1{,}536 & 2{,}776 & 1.000 & 0.941 &              & 310   & 337   \\
    \cmidrule(lr){1-10}
    Readmission  & Risk       & $\uparrow$   & 6{,}645 & 4{,}000 & 0.0001 & 0.581 & $\checkmark$ & 792   & 641   \\
                 & Confidence & $\downarrow$ & 6{,}645 & 4{,}000 & 0.0001 & 0.571 & $\checkmark$ & 792   & 641   \\
                 & Error      & $\uparrow$   & 6{,}645 & 4{,}000 & 0.0002 & 0.412 & $\checkmark$ & 792   & 641   \\
    \cmidrule(lr){1-10}
    Income       & Risk       & $\uparrow$   & 4{,}000 & 4{,}000 & 0.322 & 0.024 & $\checkmark$ & 1{,}704 & 1{,}501 \\
                 & Confidence & $\downarrow$ & 4{,}000 & 4{,}000 & 0.878 & 0.913 &              & 1{,}704 & 1{,}501 \\
                 & Error      & $\uparrow$   & 4{,}000 & 4{,}000 & 0.055 & 0.111 &              & 1{,}704 & 1{,}501 \\
    \cmidrule(lr){1-10}
    Coverage     & Risk       & $\uparrow$   & 4{,}000 & 4{,}000 & 0.0001 & 0.0001 &              & 1{,}025 & 1{,}182 \\
                  & Confidence & $\downarrow$ & 4{,}000 & 4{,}000 & 0.0001 & 0.024 &              & 1{,}025 & 1{,}182 \\
                  & Error      & $\uparrow$   & 4{,}000 & 4{,}000 & 0.0001 & 0.0001 &              & 1{,}025 & 1{,}182 \\
    \cmidrule(lr){1-10}
    NSW          & Outcome    & $\downarrow$ & 85     & 185    & 0.131 & 0.508 &              & 24     & 57    \\
    \bottomrule
  \end{tabular}
\end{table*}
